% =========================
% 1. Theorem-style problem statement
% =========================
\usepackage{amsmath,amssymb,amsthm}
\newtheorem{problem}{Problem}

\begin{problem}[Critical roughness threshold for anomalous scalar dissipation]
Consider the passive scalar advection--diffusion equation
\begin{equation}
\partial_t \theta + u\cdot \nabla \theta = D \Delta \theta,
\qquad \nabla \cdot u = 0,
\end{equation}
posed on a periodic domain in dimension $d\in\{2,3\}$, where $D>0$ is the scalar diffusivity and $u=u(x)$ is an autonomous, divergence-free random velocity field with prescribed H\"older regularity exponent $\alpha\in[-1,1]$.

We consider the following three classes of incompressible velocity fields:
\begin{enumerate}
    \item In two dimensions,
    \[
    u = \nabla^\perp \psi
    \]
    for a random stream function $\psi$.

    \item In three dimensions, a general divergence-free velocity field obtained by applying the Leray projection to a random vector field.

    \item In three dimensions, a Clebsch velocity field of the form
    \[
    u = \nabla \phi_1 \times \nabla \phi_2
    \]
    for random scalar potentials $\phi_1,\phi_2$.
\end{enumerate}

For each realization of $u$, let $\theta$ solve the scalar equation above with prescribed initial data, and define the finite-time integrated scalar dissipation
\begin{equation}
\chi(D,T)
:=
D \int_0^T \int |\nabla \theta(x,t)|^2 \, dx\,dt,
\end{equation}
where $T$ is a finite observation time chosen on the order of several large-eddy turnover times.

The problem is to determine:
\begin{enumerate}
    \item the scaling of $\chi(D,T)$ as $D\to 0$ for each of the three velocity constructions;
    \item whether anomalous dissipation occurs, i.e. whether $\chi(D,T)$ remains bounded away from zero in the vanishing-diffusivity limit after suitable ensemble averaging;
    \item the H\"older regularity exponent $\beta$ of the scalar field $\theta$ as a function of $\alpha$;
    \item whether the Yaglom-type scaling relation
    \[
    \alpha + 2\beta = 1
    \]
    holds, fails, or is only approximately valid in different roughness regimes;
    \item the critical exponent $\alpha_c$ at which anomalous dissipation first appears, and whether this threshold coincides with the onset of Yaglom-type scaling.
\end{enumerate}

A central objective is to compare these behaviors across the 2D stream-function, 3D Leray-projected, and 3D Clebsch settings, with particular emphasis on whether the more general 3D Leray class exhibits a transition distinct from that of the more structured Clebsch class.
\end{problem}



% =========================
% 2. Paper introduction paragraph
% =========================
We study the passive scalar advection--diffusion equation
\begin{equation}
\partial_t \theta + u\cdot \nabla \theta = D \Delta \theta,
\qquad \nabla\cdot u = 0,
\end{equation}
for autonomous random incompressible velocity fields in two and three spatial dimensions, with the goal of understanding the onset of anomalous scalar dissipation and its relation to scalar regularity in the vanishing-diffusivity regime. In two dimensions, the velocity field is generated by a stream function through $u=\nabla^\perp \psi$, while in three dimensions we consider both general divergence-free fields obtained by Leray projection and structured Clebsch fields of the form $u=\nabla \phi_1\times \nabla \phi_2$. For each of these three settings, we prescribe a velocity roughness exponent $\alpha\in[-1,1]$, vary the scalar diffusivity $D$ in the regime of Schmidt number less than one, and average over random realizations of the velocity field. Our primary observable is the finite-time integrated scalar dissipation
\begin{equation}
\chi(D,T)=D\int_0^T\int |\nabla \theta(x,t)|^2\,dx\,dt,
\end{equation}
measured over finite but dynamically relevant times $T$, since as $t\to\infty$ the scalar relaxes to its spatial mean. We also estimate the H\"older exponent $\beta$ of the scalar field and test the Yaglom-type prediction $\alpha+2\beta=1$. The central question is whether there exists a critical roughness threshold $\alpha_c$ separating a regime without anomalous dissipation from one in which dissipation persists as $D\to 0$, and whether this transition coincides with the recovery of Yaglom-type scaling. A key motivation is to compare the extent to which this transition depends on the geometric structure of the incompressible velocity field, especially in the largely open three-dimensional Leray-projected setting.




% =========================
% 3. Grant-style research objective
% =========================
\subsection*{Research Objective}

The objective of this project is to determine the critical velocity roughness threshold governing the onset of anomalous dissipation for passive scalars advected by autonomous random incompressible flows in two and three dimensions. We consider the advection--diffusion equation
\begin{equation}
\partial_t \theta + u\cdot \nabla \theta = D\Delta \theta,
\qquad \nabla\cdot u = 0,
\end{equation}
and investigate three distinct classes of divergence-free velocity fields: (i) two-dimensional stream-function-generated flows $u=\nabla^\perp\psi$, (ii) three-dimensional flows obtained by Leray projection, and (iii) three-dimensional Clebsch flows $u=\nabla\phi_1\times\nabla\phi_2$. Across these settings, we prescribe a broad range of velocity H\"older exponents $\alpha$, from smooth to nearly critical rough fields, and we vary the scalar diffusivity $D$ in the regime $Sc<1$.

Our first aim is to quantify the vanishing-diffusivity scaling of the finite-time integrated scalar dissipation
\begin{equation}
\chi(D,T)
=
D\int_0^T\int |\nabla\theta(x,t)|^2\,dx\,dt,
\end{equation}
where $T$ is chosen to be a finite observation window on the order of several large-eddy turnover times. This finite-time perspective is essential, since at asymptotically long times the scalar field relaxes to its mean and dissipation measurements become trivial. Our second aim is to estimate the H\"older regularity exponent $\beta$ of the scalar field and to test the validity of the Yaglom-type relation
\begin{equation}
\alpha + 2\beta = 1.
\end{equation}
Together, these measurements will identify whether anomalous dissipation and Yaglom-type scaling emerge simultaneously or at distinct thresholds.

The central hypothesis of the project is that there exists a critical exponent $\alpha_c$ separating a non-anomalous regime from an anomalous one, and that this threshold depends on the structural class of the velocity field. In particular, we seek to determine whether the transition observed for two-dimensional and Clebsch-generated flows persists for general three-dimensional Leray-projected flows, where the role of geometric constraints is much less understood. By systematically varying $\alpha$, $D$, and the random seed used to generate the velocity field, and by ensemble-averaging the resulting statistics, this project will provide a quantitative map of the dissipation and regularity regimes for passive scalar transport in rough incompressible flows.






% =========================
% 4. concise
% =========================


\begin{problem}
For the passive scalar equation
\[
\partial_t\theta + u\cdot\nabla\theta = D\Delta\theta,
\qquad \nabla\cdot u=0,
\]
with autonomous random divergence-free velocity field $u$ of H\"older regularity exponent $\alpha\in[-1,1]$, determine the critical exponent $\alpha_c$ at which anomalous scalar dissipation first occurs as $D\to 0$. This question is to be investigated in three settings: (i) $d=2$ with $u=\nabla^\perp\psi$, (ii) $d=3$ with $u$ obtained by Leray projection, and (iii) $d=3$ with Clebsch velocity field $u=\nabla\phi_1\times\nabla\phi_2$.

For each setting, analyze the scaling as $D\to0$ of
\[
\chi(D,T)=D\int_0^T\int |\nabla\theta(x,t)|^2\,dx\,dt,
\]
for finite observation times $T$ of several large-eddy turnover times, and estimate the scalar H\"older exponent $\beta$. The goal is to identify the dependence of anomalous dissipation on $\alpha$, to determine whether the Yaglom-type relation
\[
\alpha+2\beta=1
\]
holds in any asymptotic regime, and to compare the transition mechanisms across the three velocity constructions.
\end{problem}